\section{Introduction}
Human activity recognition (HAR) from smartphone inertial sensors is a standard pattern
recognition task: accelerometer and gyroscope windows are summarized into hand-crafted
time- and frequency-domain features and classified into activities~\cite{anguita2013}. In
deployment, however, the phone that produced the training data is rarely the phone that is
used. Phones differ in sensor bias, sensitivity and noise, and models trained on one device
degrade on others~\cite{stisen2015}. In the dataset-shift vocabulary this is a covariate
shift between a source and a target domain~\cite{quinonero2009,moreno2012}; theory bounds
the target error by the source error plus a divergence between the domains~\cite{bendavid2010},
so reducing that divergence matters as much as a strong classifier.

\textbf{Problem.} Given records of 561 features from Device\_B (\NTrain{} labelled records),
predict the activity of records from unseen devices. The validation set (\NVal{} records)
turned out to contain two devices; the hidden test set may contain others.

\textbf{Related work.} Domain generalization (DG) learns from source domains only, while
domain adaptation also uses unlabeled target data~\cite{zhou2022dg,wang2022dg}. A simple and
effective family re-estimates normalization statistics on the target domain:
AdaBN~\cite{li2017adabn} recomputes batch-normalization statistics per domain, and
CORAL~\cite{sun2016coral} aligns second-order statistics. Per-device calibration and
normalization is also the first mitigation proposed for sensing heterogeneity
in~\cite{stisen2015}. Deep models dominate HAR on raw signals~\cite{wang2019deep}, but on
hand-crafted feature vectors kernel and linear models remain competitive~\cite{anguita2013}.
% TODO(team, T-014): verify every citation; add 2-4 HAR domain-adaptation papers you read.

\textbf{Contributions.} (i) A diagnosis of the released data: per-feature affine device shift,
two devices in the validation set, and a noise-dominated feature block. (ii) An unsupervised
device-discovery step on activity-independent probe features that enables per-device
standardization at test time. (iii) A complete, tested, sub-millisecond pipeline with an
honest timing protocol, evaluated by in-domain CV, simulated-shift CV and validation.
